% !TeX root = ../../Book3.tex
% 纲要标题：## 附录 E　微分代数与微分消去

\chapter{微分代数与微分消去}
\label{b3:app:E}
\BookChapterNumber{b3:app:E}

% 纲要细目（与计划纲要同步）：
% > 本附录简要介绍微分代数。主体采用特征零单导子微分域；多交换导子基定理另外注明来源与范围。所需基础是域、理想、局部化和多项式；交换代数正文不依赖本附录。非线性算法的完整证明和实现留给专题教材。

\begin{chapterabstract}
本附录建立导子、微分多项式与微分理想的基本运算，比较根微分理想的有限性和普通 Noether 性。特征集与 Rosenfeld--Gröbner 方法处理消去中的非零及退化分支；线性微分模与 Picard--Vessiot 扩张则给出另一类结构问题。较深的有限性与分解结果只给证明概要。
\end{chapterabstract}

\begin{learninggoals}
\begin{itemize}
\item 给一个环指定导子，并在商环、分式域与微分多项式环上检查相容性。
\item 区分普通代数次数、微分阶及两种超越次数。
\item 在微分消去中保存乘子证书，记录初式、分离子与非零条件，核验退化分支。
\item 从线性微分系统导出常数矩阵作用，并辨认其与完整微分 Galois 理论的差别。
\end{itemize}
\end{learninggoals}

把 \(y,y',y'',\ldots\) 暂时当作独立变量，若只施加关系 \(y'-y=0\)，普通理想商允许用 \(y\) 代替 \(y'\)，却未必允许把求导作用降到商上：这个关系的导数 \(y''-y'\) 还没有被置零。要在商中继续微分，必须连同原关系的所有后续导数一起处理。微分理想由此成为微分方程的合适代数载体。

\ProseParagraphBreak

一旦把所有导数加入，变量和可能出现的关系都变成无限族；再做消去时，还要注意某个系数为零会使方程阶数突然下降。特征集方法需要把非零系数分支与退化分支分开，才能既不漏解也不错误地除以零。线性方程则可由微分模组织，形成与非线性消去不同的结构问题。

\ProseParagraphBreak

所需基础为域、理想、多项式与局部化，Kähler 微分的比较用到本卷第14章。本附录的特征集与算法说明采用一个导子、特征零；多个两两交换导子的基定理另行注明来源。普通常微分方程的解析存在定理不作为预备知识，因为这里的定义与多数算例可以在形式代数中完成。Hubert 的正式文章 \textit{Notes on Triangular Sets and Triangulation-Decomposition Algorithms II: Differential Systems} 的 Sections 2--3 对应基本对象与特征集\cite[Sections 2--3]{Hubert2003DifferentialSystems}；完整 Rosenfeld--Gröbner 输出规格见 Boulier 等的 Section 6.2\cite[Section 6.2, Theorem 60]{BoulierLazardOllivierPetitot2009DifferentialIdeals}。

\ProseParagraphBreak
前两节建立导子、微分理想及其几何语言；根微分理想的有限性证明还要用第三节的特征集性质。第三节说明排序、饱和分支与计算证书，Rosenfeld--Gröbner 的完整算法及终止技术采用文献。第四节直接建立线性微分模、PV 条件与常数矩阵作用，随后介绍微分闭域、参数化理论与差分结构；这些后续对象的存在性、一般对应及模型论结果没有在这里系统展开。
\begin{center}
\begin{tabularx}{\linewidth}{@{}>{\raggedright\arraybackslash}p{0.38\linewidth}X@{}}
\toprule
所证明的代数结论 & 仍需另作论证的解析问题 \\
\midrule
微分多项式关系及理想成员 & 指定形式解的收敛与函数定义域 \\
微分扩域中的一个解 & 给定初值的局部解析存在与唯一性 \\
带非零条件的微分消去 & 分母零点或奇点处的解析延拓 \\
PV 扩张及其自同构 & 一般常微分方程的解析分类 \\
\bottomrule
\end{tabularx}
\end{center}

% !TeX root = ../../Book3.tex
% 纲要标题：### E.1 微分结构与基本对象

\section{微分结构与基本对象}
\label{b3:sec:E01}

普通多项式关系只约束元素本身，微分方程还把元素与反复求导后的值联系起来。把导子作为环的指定结构后，可以对这些关系作理想运算；常数域也随所选导子而确定。

% 纲要细目（与计划纲要同步）：
% * 指定导子、Leibniz 法则、常数、微分理想、商环及分式域上的延拓；分式域延拓适用于任意整环
% * \(K\{y_1,\ldots,y_n\}=K[y_i^{(j)}:j\ge0]\)、微分阶及微分域扩张；指数函数作为微分代数元的例子
% * 区别指定导子的微分代数、Kähler 微分和 DG 代数，比较值域、Leibniz 法则与平方零条件
% #### E.1.1 导子、常数、微分理想与分式域延拓
% #### E.1.2 微分多项式、微分阶与微分域扩张
% #### E.1.3 微分代数、Kähler 微分与 DG 代数

\subsection{导子、常数、微分理想与分式域延拓}
\label{b3:subsec:E01:01}

除另有说明外，本附录的域均为特征零域，且只指定一个导子 \(\delta\)。例如方程 \(\delta(y)=y\) 同时约束元素与它的导数；取商时，也必须使这个指定导子在商环上有定义。

\begin{definition}\label{b3:def:differential-ring}
设 \(A\) 为交换含幺环，\(\delta:A\to A\) 为加法映射。若对任意 \(a,b\in A\) 都满足
\[
\delta(ab)=a\cdot\delta(b)+\delta(a)\cdot b.
\]
则称 \(\delta\) 为 \(A\) 的\term{导子}[derivation]，称 \((A,\delta)\) 为\term{微分环}[differential ring]。若 \(A\) 是域，则称 \((A,\delta)\) 为微分域。称 \(C_A=\ker\delta\) 为常数环；若 \(A\) 的理想 \(I\) 满足 \(\delta(I)\subseteq I\)，则称 \(I\) 为\term{微分理想}[differential ideal]。
\end{definition}

Leibniz 法则给出 \(\delta(1)=0\)，也给出 \(C_A\) 对加法与乘法封闭。在微分域中，\(\delta(a^{-1})=-a^{-2}\cdot\delta(a)\)，所以非零常数的逆仍为常数，\(C_A\) 是域。微分同态是与导子交换的环同态；其核为微分理想。反之，\(\delta(\bar a)=\overline{\delta(a)}\) 在 \(A/I\) 上良定义，恰好要求 \(I\) 是微分理想。

\begin{proposition}\label{b3:prop:derivation-fraction-extension}
设 \(A\) 为整环，\(\delta:A\to A\) 为导子。它唯一延拓到 \(K=\operatorname{Frac}(A)\)，且对任意 \(a,b\in A\)、\(b\ne0\)，有
\[
\delta(a/b)=\frac{b\cdot\delta(a)-a\cdot\delta(b)}{b^2}.
\]
\end{proposition}
\begin{proof}
在域中对 \(bb^{-1}=1\) 求导强制给出上述公式，故至多有一个延拓。若 \(a/b=c/d\)，则 \(ad=bc\)。对此等式求导，再用 \(ad=bc\) 整理，得到
\[
d^2\bigl(b\cdot\delta(a)-a\cdot\delta(b)\bigr)
=b^2\bigl(d\cdot\delta(c)-c\cdot\delta(d)\bigr).
\]
故公式不依赖代表元。通分直接验证加法与 Leibniz 法则，于是确实得到导子。
\end{proof}

例如在 \(\mathbb Q(t)\) 上指定 \(\delta(t)=1\)，即通常的形式求导。若改为零导子，同一个底层域便成为另一微分域，其常数域也随之改变。因此仅给出一个环，还没有指定微分结构。

\subsection{微分多项式、微分阶与微分域扩张}
\label{b3:subsec:E01:02}

\begin{definition}\label{b3:def:differential-polynomial-ring}
设 \((K,\delta)\) 为特征零微分域，\(n\ge1\)。取彼此代数独立的符号 \(y_i^{(j)}\)（\(1\le i\le n,j\ge0\)），在多项式环
\[
K\{y_1,\ldots,y_n\}=K[y_i^{(j)}:1\le i\le n,\ j\ge0]
\]
上把导子延拓为 \(\delta(y_i^{(j)})=y_i^{(j+1)}\)。所得环称为\term{微分多项式环}[differential polynomial ring]。记 \(y_i^{(0)}=y_i\)。非恒定微分多项式的\term{微分阶}[differential order]是其中实际出现的最高导数阶数。
\end{definition}

每个微分多项式仍只有有限项。例如 \((y')^2-y\) 为一阶、通常多项式次数二，而 \(y^{(3)}-t y\) 为三阶、关于全部导数变量的次数一。次数与微分阶分别记录非线性程度和所用导数深度。

\begin{proposition}\label{b3:prop:differential-polynomial-universal}
设 \((K,\delta)\) 为特征零微分域，\(L/K\) 为保持导子的微分域扩张，\(a_1,\ldots,a_n\in L\)。存在唯一微分 \(K\) 代数同态
\[
\operatorname{ev}_a:K\{y_1,\ldots,y_n\}\longrightarrow L,
\qquad y_i^{(j)}\longmapsto\delta^j(a_i).
\]
\end{proposition}
\begin{proof}
多项式环的泛性质给出唯一普通代数同态。它在系数与所有生成元上和导子交换，再由加法及 Leibniz 法则知在每个多项式上交换。反过来，任一微分同态的所有生成元像都已被 \(y_i\mapsto a_i\) 强制决定。
\end{proof}

设 \((K,\delta)\) 为特征零微分域，\(L/K\) 为保持导子的微分域扩张，\(a\in L\)。若存在非零 \(P\in K\{y\}\) 使 \(P(a)=0\)，则称 \(a\) 为 \(K\) 上的\term{微分代数元}[differentially algebraic element]；否则称为 \(K\) 上的微分超越元。这里 \(P(a)\) 把 \(y^{(j)}\) 替换为 \(\delta^j(a)\)。在带通常求导的亚纯函数域中，\(e^t\) 满足 \(y'-y=0\)，所以是 \(\mathbb C(t)\) 上的微分代数元。普通代数性要求零阶关系，微分代数性则允许导数参与。

\ProseParagraphBreak
记 \(K\langle a\rangle=K(a,\delta a,\delta^2a,\ldots)\)，这是包含 \(K,a\) 的最小微分子域。若 \(a\) 满足关系，并不表示其任意若干初值都能实现为解析解；代数扩张与解析存在性须分别讨论。

\subsection{微分代数、Kähler 微分与 DG 代数}
\label{b3:subsec:E01:03}

设 \(k\) 为域，\(A\) 为交换 \(k\) 代数。本卷第14章的 Kähler 微分 \(\mathrm{d}:A\to\Omega_{A/k}\) 同时表示所有 \(k\) 导子，而这里的 \(\delta:A\to A\) 是一个已经指定的导子。如果 \(\delta|_k=0\)，泛性质给出唯一 \(A\) 线性映射 \(u:\Omega_{A/k}\to A\) 使 \(\delta=u\circ \mathrm{d}\)。因此微分模并不自动选出 \(\delta\)。例如 \(A=k[x]\) 时，\(\Omega_{A/k}=A\,\mathrm{d}x\)，选定 \(u(\mathrm{d}x)=h(x)\) 才得到导子
\[
\delta=h(x)\frac{\mathrm{d}}{\mathrm{d}x}.
\]

另一个不同对象是\term{微分分次代数}[differential graded algebra, DG algebra]。设 \(B=\bigoplus_iB^i\) 为分次 \(k\) 代数，\(\mathrm{d}:B^i\to B^{i+1}\) 为次数一的 \(k\) 线性映射。若 \(\mathrm{d}^2=0\)，且对任意齐次元素 \(a,b\in B\) 满足
\[
\mathrm{d}(ab)=\mathrm{d}(a)b+(-1)^{\deg a}a\,\mathrm{d}(b),
\]
则称 \((B,\mathrm{d})\) 为微分分次代数。本附录的普通导子不要求平方为零。例如在特征零域 \(K\) 的多项式环 \(K[t]\) 中，\(\delta=\mathrm{d}/\mathrm{d}t\) 满足 \(\delta^2(t^2)=2\ne0\)。\cref{b3:tab:E-three-differentials}比较了三种结构：指定导子选择一种求导运算，Kähler 微分表示全部导子，DG 微分则还带有分次和平方零条件。

\begin{table}[htbp]
\centering
\caption{指定导子、Kähler 微分与 DG 微分的比较}
\label{b3:tab:E-three-differentials}
\begin{tabularx}{\linewidth}{@{}>{\raggedright\arraybackslash}p{0.20\linewidth}>{\raggedright\arraybackslash}p{0.28\linewidth}>{\raggedright\arraybackslash}X>{\raggedright\arraybackslash}p{0.17\linewidth}@{}}
\toprule
结构 & 映射与值域 & Leibniz 法则 & 平方零条件 \\
\midrule
指定导子 & \(\delta:A\to A\) & 普通乘积法则 & 不要求 \(\delta^2=0\) \\
Kähler 微分 & \(\mathrm{d}:A\to\Omega_{A/k}\) & 普通乘积法则 & 此映射不能自复合 \\
DG 微分 & \(\mathrm{d}:B^i\to B^{i+1}\) & 含 \((-1)^{\deg a}\) 的分次乘积法则 & 要求 \(\mathrm{d}^2=0\) \\
\bottomrule
\end{tabularx}
\end{table}

多导子版本指定有限个 \(\delta_1,\ldots,\delta_m\)，并要求它们两两交换；生成元变成 \(\delta_1^{\alpha_1}\cdots\delta_m^{\alpha_m}y_i\)。交换性对应混合偏导次序相同。普通微分情形只有一个导子，偏微分情形则要同时处理各个多重导数之间的相容性。基本定义及其计算记号可参阅\cite[Section 3]{BoulierLazardOllivierPetitot2009DifferentialIdeals}。

% !TeX root = ../../Book3.tex
% 纲要标题：### E.2 有限性与微分代数几何的基本概念

\section{有限性与微分代数几何的基本概念}
\label{b3:sec:E02}

微分多项式环含有无限多个导数变量，普通 Noether 性因而失效。特征零下，根微分理想仍有有限性与分解定理；这些结论约束的是保持求导的理想，不能直接套用到任意普通理想。

% 纲要细目（与计划纲要同步）：
% * 普通环的 Noether 性与 Ritt–Raudenbush 根微分理想升链条件不同；注明特征零、有限多个交换导子和有限多个微分不定元
% * 微分代数集、Kolchin 拓扑与微分超越次数；明确工作微分域和所取扩域
% * 用一组低阶方程说明代数关系、微分关系与常数域变化；代数相容性不替代解析存在、收敛和初值唯一性
%
% #### E.2.1 Noether 性与 Ritt–Raudenbush 升链条件
% #### E.2.2 微分代数集、Kolchin 拓扑与微分超越次数
% #### E.2.3 低阶微分方程、常数域与解析解


% 正文已按纲要写入；有限性与算法长证明给出概要及出处。

\subsection{Noether 性与 Ritt–Raudenbush 升链条件}
\label{b3:subsec:E02:01}

虽然每个微分多项式只用有限多个变量，整个环 \(K\{y\}\) 却不是 Noether 环：普通理想形成严格升链
\[
(y)\subsetneq(y,y')\subsetneq(y,y',y'')\subsetneq\cdots.
\]
它们一般不是微分理想。例如 \(y'\notin(y)\)。由一个集合 \(F\) 生成的最小微分理想记为 \([F]\)，具体就是所有 \(\delta^jf\)（\(f\in F,j\ge0\)）生成的普通理想。

\begin{lemma}\label{b3:lem:differential-radical}
设 \(A\) 为含 \(\mathbb Q\) 的微分环，\(I\) 为微分理想。则普通根理想 \(\sqrt I\) 也是微分理想。
\end{lemma}
\begin{proof}
导子下降到 \(A/I\)，故只须证明导子保持幂零元。若 \(a^n=0\)，反复用 Leibniz 法则展开 \(\delta^n(a^n)=0\)。其中对每个因子恰求导一次的项为 \(n!(\delta a)^n\)，其他项至少有一个因子未经求导，因此属于 \((a)\)。因 \(n!\) 可逆，得 \((\delta a)^n\in(a)\)，再取 \(n\) 次幂得到 \((\delta a)^{n^2}=0\)。
\end{proof}

\begin{theorem}[Ritt--Raudenbush 定理]\label{b3:thm:ritt-raudenbush}
设 \(K\) 为特征零微分域，指定有限个两两交换的导子，并取有限个微分不定元。则微分多项式环的根微分理想满足升链条件。等价地，每个根微分理想均形如 \(\sqrt{[f_1,\ldots,f_r]}\)，其中 \(r\) 有限。
\end{theorem}

\begin{proofsketch}
固定一个与各导子相容的导数排序。\secref{b3:sec:E03}定义互约化集及特征集：按导变量及其次数从小到大比较，同一前缀时较长列表为较小秩。Dickson 引理保证互约化集有限，且这些列表的秩不存在无限严格下降；微分伪约化又给出
\(hP-r\in[\mathcal A]\)，其中 \(r\) 对 \(\mathcal A\) 约化，\(h\) 为初式与分离子的幂乘积。这里省略列表秩的良基性与伪约化的完整归纳，构造及证明见\cite[Sections 3.1--3.3]{Hubert2003DifferentialSystems}。

设 \(J_1\subseteq J_2\subseteq\cdots\) 为根微分理想链，\(J=\bigcup_iJ_i\)。若 \(J=0\) 或 \(J\) 为全环，结论立即成立。否则取 \(J\) 的一个特征集 \(\mathcal A\)。它有限，故包含在某个 \(J_{i_0}\) 中；对 \(i\ge i_0\)，其最小秩性质使它同时成为 \(J_i\) 的特征集。记 \(h_1,\ldots,h_t\) 为 \(\mathcal A\) 的初式与分离子，\(H=\prod_jh_j\)。每个 \(P\in J_i\) 的约化余式也属于 \(J_i\)，若非零便会产生更小秩的互约化集，故余式为零。于是
\[
J_i:H^\infty=[\mathcal A]:H^\infty\qquad(i\ge i_0),
\]
非零分支已经稳定。

根微分理想满足分支恒等式
\[
J_i=(J_i:H^\infty)\cap
\bigcap_{j=1}^t\sqrt{[J_i,h_j]}.
\]
它可由极小素理想检查：特征零时，根微分理想的极小素理想仍受各导子保持，局部化到该极小素点所得约化零维环是域，便可核验这一性质；每个这样的素理想或者不含 \(H\)，或者含某个 \(h_j\)。相应的饱和与添入方程分支恰好覆盖这两种情况。

各非零 \(h_j\) 都对 \(\mathcal A\) 约化且秩较低；将它添入理想后，特征集秩严格降低，或直接得到全环。因此对特征集秩作良基归纳，各链 \(\sqrt{[J_i,h_j]}\) 都稳定。分支只有有限个，结合已经稳定的饱和部分，原链也稳定。这是特征集证明的主要归纳步骤；准确结论见\cite[Section 2.2, Theorem 2.4]{Hubert2003DifferentialSystems}。

最后，升链条件使逐次选取新根微分生成元的过程终止。反之，对链的并取有限个根微分生成元，它们已同时落在某一项中，该项以后便不再增长。于是两种表述等价。
\end{proofsketch}

结论只针对根微分理想，不能改写成整个微分多项式环是 Noether 环；特征零还用于分离子非零及根理想受导子保持。

\subsection{微分代数集、Kolchin 拓扑与微分超越次数}
\label{b3:subsec:E02:02}

在一个微分扩域 \(L\supseteq K\) 中，对 \(F\subseteq K\{y_1,\ldots,y_n\}\) 定义
\[
V_L(F)=\{a\in L^n:P(a)=0\text{ 对所有 }P\in F\}.
\]
这些集合称为\term{微分代数集}[differential algebraic set]。任意交由合并方程定义，有限并由两组方程的所有两两乘积定义，因为 \(L\) 是域；故它们构成某个拓扑的闭集，这就是\term{Kolchin 拓扑}[Kolchin topology]。

\ProseParagraphBreak
\(V_L(F)=V_L(\sqrt{[F]})\)，因为微分求值保持导子，而域无非零幂零元。逆向恢复理想时，需要充分多的解；任意给定的 \(L\) 未必有这种性质。例如在 \(K=\mathbb C\) 上指定零导子，方程 \(y'=1\) 在 \(K\) 内无解，却在 \(\mathbb C(t)\) 中有解 \(t\)。不能把“在当前底域没有解”误当成微分理想为全环。

\begin{definition}\label{b3:def:differential-transcendence}
设 \((L,\delta)/(K,\delta)\) 为带同一导子的微分域扩张，\((a_i)_{i\in I}\) 为 \(L\) 中的元素族。若全部 \(\delta^ja_i\)（\(i\in I,j\ge0\)）在 \(K\) 上代数独立，则称该族\term{微分代数独立}[differentially algebraically independent]。\(L/K\) 的\term{微分超越次数}[differential transcendence degree]定义为这类元素族的基数的上确界。
\end{definition}

这与普通超越次数不同。在带 \(\delta t=1\) 的 \(K(t)\) 中，\(t\) 满足 \(t'-1=0\)，故微分超越次数为零，普通超越次数却为一。另一方面，\(K\langle y\rangle\) 中各 \(y^{(j)}\) 都独立，它有一个微分超越生成元，但普通超越次数无限。基的交换性质与一般维数理论需要进一步论证；本附录的算例均直接从生成元的关系判断。

\subsection{低阶微分方程、常数域与解析解}
\label{b3:subsec:E02:03}

微分方程中的“常数”首先指导子为零的元素。对 \(y'=0\)，解集正是所在扩域的常数域；换一个常数更大的扩域就会出现更多解。对 \(y'=y\)，若 \(u,v\) 是两个非零解，则
\[
\delta(u/v)=\frac{u'v-uv'}{v^2}=0,
\]
所以它们相差一个非零常数倍。这个结论纯属域中的运算，无须先定义指数函数。

\begin{example}\label{b3:exm:formal-riccati}
令 \(K\) 为特征零域，\(K[\![t]\!]\) 上取逐项导子。对给定 \(a\in K\)，方程 \(y'=y^2,y(0)=a\) 有唯一形式幂级数解
\[
y=\frac{a}{1-at}=\sum_{n\ge0}a^{n+1}t^n.
\]
确实，若 \(y=\sum_nc_nt^n\)，系数比较给出 \((n+1)c_{n+1}=\sum_{i=0}^nc_ic_{n-i}\)，由 \(c_0=a\) 递归唯一决定全部系数；所写几何级数直接满足等式。
\end{example}

当 \(K=\mathbb C\) 时，上例还在足够小邻域内收敛；这是该具体几何级数的性质。一般形式解的收敛、实区间上的存在时长、经过奇点的延拓，都属于额外的解析问题。微分代数能证明一个方程由其他方程推出，却不会仅凭理想运算自动提供这些解析结论。

% !TeX root = ../../Book3.tex
% 纲要标题：### E.3 特征集与有效消去的代表方法

\section{特征集与有效消去的代表方法}
\label{b3:sec:E03}

消去导数变量时，最高导数的选择决定约化次序，初式与分离子则决定哪些步骤可以除去因子。本节在带一个导子的特征零微分域上讨论：先把这种选择写成排序与特征集，再说明局部化后被倒置因子的零点分支为何仍须处理。

% 纲要细目（与计划纲要同步）：
% * 导数排序、列表秩、自约化、初式、分离子及伪约化；用正式非首一算例完整计算导变量、初式、分离子和导数
% * 在特征集之后给出 Ritt–Raudenbush 证明概要，说明主饱和分支稳定与低秩退化归纳；技术引理注明原文来源
% * 比较正则微分系统、正则微分链与特征展示的不同余式判据；一般 Rosenfeld–Gröbner 更新、正确性与终止性保留文献输入
% * 完整证明乘子证书与退化分支算例；原相关练习改为排序比较和非齐次消去，区别清分母所给的必要条件与等价系统
% #### E.3.1 导数排序、导变量、初式、分离子与伪约化
% #### E.3.2 Rosenfeld–Gröbner 正则微分链与饱和分解
% #### E.3.3 消去算例与分母、初式、分离子的退化分支

\subsection{导数排序、导变量、初式、分离子与伪约化}
\label{b3:subsec:E03:01}

要把求导后的方程组织成可消去的形式，必须先决定优先消去哪一个导数。导数排序首先比较的是导数变量；第8章的单项式序则比较整个单项式，二者不能直接等同。

\begin{definition}\label{b3:def:differential-ranking}
设 \((K,\delta)\) 为带一个导子的特征零微分域，\(R=K\{y_1,\ldots,y_n\}\)，\(n\ge1\)，全部导数变量组成集合
\(\mathcal U=\{\delta^jy_i:1\le i\le n,j\ge0\}\)。若 \(\mathcal U\) 上的全序满足
\[
u<\delta u,\qquad u<v\Longrightarrow\delta u<\delta v,
\]
则称其为\term{导数排序}[differential ranking]。
\end{definition}

例如 \(y<z<y'<z'<\cdots\) 按阶优先；\(y<y'<\cdots<z<z'<\cdots\) 则优先消去 \(z\) 的导数。这两种排序都是良序。一般单导子排序也是良序：若存在无限严格下降序列，有限个微分不定元中必有某个 \(y_i\) 的导数出现无限多次。取这些项形成的子列，它仍严格下降；但 \(y_i<\delta y_i<\delta^2y_i<\cdots\)，所以对应导数阶数须形成非负整数的无限严格下降序列，不可能。

\begin{definition}\label{b3:def:differential-leader}
设 \((K,\delta)\) 为特征零微分域，\(R=K\{y_1,\ldots,y_n\}\)，并固定导数排序。若 \(P\in R\setminus K\) 出现的最高导数变量为 \(u\)，把它写成
\[
P=I_Pu^d+\sum_{j=0}^{d-1}a_ju^j,\qquad d\ge1,
\]
其中各系数只含低于 \(u\) 的变量，则称 \(u\) 为\term{导变量}[leader]，\(I_P\) 为\term{初式}[initial]，\(S_P=\partial P/\partial u\) 为\term{分离子}[separant]。定义多项式的秩为
\[
\operatorname{rank}P=(u,d),
\]
先按导变量排序比较，再在导变量相同时比较正整数次数。
\end{definition}
\symidx{\(I_P\)、\(S_P\)：微分多项式 \(P\) 的初式与分离子}
\symidx{\(\operatorname{rank}P\)：微分多项式 \(P\) 的秩}

由于 \(K\) 的特征零，\(S_P\ne0\)。求导后最高新变量为 \(\delta u\)，并有 \(\delta P=S_P\delta u+\text{较低变量的项}\)。因此消去 \(u\) 时可能需要乘初式，消去更高导数时需要乘分离子。例如 \(P=(y')^2-y\) 的初式为一，分离子为 \(2y'\)，而 \(\delta P=y'(2y''-1)\)。在 \(y'\ne0\) 的分支可推出 \(2y''-1=0\)；直接约去 \(y'\) 则会遗漏 \(y=0\) 这一解。

非首一的例子能同时显示两类因子：对 \(P=y(y'')^3+(y')^2-1\)，导变量为 \(y''\)，初式为 \(y\)，分离子为 \(3y(y'')^2\)。初式为零会使关于导变量的次数下降，分离子为零则会使求导后新变量的系数为零；完整求导计算见本节末的\cref{b3:exm:E-initial-separant}。

\begin{definition}\label{b3:def:differential-characteristic-set}
设 \((K,\delta)\) 为特征零微分域，\(R=K\{y_1,\ldots,y_n\}\)，并固定导数排序。对导变量为 \(u\)、关于 \(u\) 次数为 \(d\) 的 \(P\in R\setminus K\)，若 \(Q\in R\) 不含 \(u\) 的真导数 \(\delta^ju\)（\(j\ge1\)），则称 \(Q\) 对 \(P\) 部分约化；若还满足 \(\deg_uQ<d\)，则称它对 \(P\) 完全约化。一个集合的每个成员若都对其余成员完全约化，则称该集合\term{自约化}[autoreduced]。

把有限自约化集合按秩递增写成列表 \(\mathcal A=(A_1,\ldots,A_s)\)。列表秩比较采用字典序：首个不同位置的多项式秩较小者，列表秩较小；若一条列表是另一条的真前缀，则较长的列表秩较小。空列表的秩因而大于任意非空列表。若 \(J\subsetneq R\) 是微分理想，\(\mathcal A\subseteq J\setminus K\) 是其中秩最小的自约化集合，则称其为 \(J\) 的\term{特征集}[characteristic set]；零理想的特征集取空集。
\end{definition}

真前缀情形把较长列表规定为较低秩，是 Ritt 的列表秩约定：前面各秩相同时，能再加入一个约化方程表示掌握了更多约束，这不是通常有限字典序的长度约定。若两个导变量为 \(\delta^ry_i\)、\(\delta^sy_i\)，其中 \(r<s\)，后者是前者的真导数，高阶成员就不可能对低阶成员部分约化。因此每个 \(y_i\) 至多贡献一个导变量，自约化集合的长度至多为 \(n\)。多项式秩是良序，而列表长度有共同上界，故上述列表秩没有无限下降链，特征集由此存在。特征集的极小性还意味着：理想中对它完全约化的多项式只能为零。若出现非零约化多项式，把它插入低于其秩的前缀并舍去后续项，便得到秩更小的自约化列表；详细的秩比较见\cite[Section 3.3, Proposition 3.15 and Definition 3.17]{Hubert2003DifferentialSystems}。

\ProseParagraphBreak
对一个有限自约化集合 \(\mathcal A\)，微分伪约化逐次使用其成员及导数，输出完全约化的余式 \(r\)，并保存
\[
hP-r\in[\mathcal A],
\]
其中 \(h\) 是 \(\mathcal A\) 的初式与分离子的幂乘积。消去导变量的真导数时，使用相应方程的导数及其分离子；消去导变量本身的高次幂时，使用初式作普通伪除法。完整步骤及终止论证见\cite[Section 3.1]{BoulierLazardOllivierPetitot2009DifferentialIdeals}。特征集通常不是原理想的普通生成集；对任意自约化集合，余式为零直接证明的是某个 \(hP\) 落在 \([\mathcal A]\) 中，不能无条件去掉 \(h\)。

乘子 \(h\) 记录实际被倒置的初式与分离子，是附录 B 所强调的有限证书的一部分。若 \(H\) 是这些因子的乘积，\(r=0\) 立即给出 \(P\in[\mathcal A]:H^\infty\)，而非无条件的 \(P\in[\mathcal A]\)。两方程的具体乘子等式及零分支核验见本节末的\cref{b3:exm:E-branch-certificate}。

下面的有限性论证先对链的并取特征集，等其有限成员全部进入某一项后稳定非零饱和分支；退化分支由更低秩的良基归纳稳定，再用已证明的分支恒等式合回原链。特征集的秩比较和伪约化性质仍采用所引文献，下面展开的是这些输入如何推出升链条件。

下面给出\cref{b3:thm:ritt-raudenbush}的证明概要。
\begin{proofsketch}
设 \((K,\delta)\) 为特征零微分域，\(R=K\{y_1,\ldots,y_n\}\)，并固定上面定义的导数排序。采用两个特征集性质：自约化列表的秩构成良基次序；若 \(\mathcal A\) 是微分理想 \(J\) 的特征集，则 \(J\) 中对 \(\mathcal A\) 完全约化的多项式只能为零。伪约化保存
\(hP-r\in[\mathcal A]\)，其中 \(h\) 为初式与分离子的幂乘积，\(r\) 完全约化。这些约化性质的证明见\cite[Sections 3.2--3.3]{Hubert2003DifferentialSystems}。

对特征集秩作良基归纳。设 \(J_1\subseteq J_2\subseteq\cdots\) 为根微分理想链，\(J=\bigcup_iJ_i\)。若 \(J=0\)，则每项均为零；若 \(J=R\)，则 \(1\) 已属于某一项，链从该项起稳定。其余情形取 \(J\) 的一个特征集 \(\mathcal A\)。它有限，故包含在某个 \(J_{i_0}\) 中。对 \(i\ge i_0\)，\(\mathcal A\) 同时是 \(J_i\) 的特征集：否则 \(J_i\subseteq J\) 内已有秩更小的自约化集，违背所选特征集的极小性。

记 \(h_1,\ldots,h_t\) 为 \(\mathcal A\) 的初式与分离子，\(H=\prod_jh_j\)。任意 \(P\in J_i\) 的伪余式 \(r=hP-\sum_jq_j\delta^{e_j}A_j\) 仍属于 \(J_i\)，所以 \(r=0\)。于是每个 \(P\in J_i\) 都属于 \([\mathcal A]:H^\infty\)。结合 \([\mathcal A]\subseteq J_i\) 得
\[
J_i:H^\infty=[\mathcal A]:H^\infty\qquad(i\ge i_0),
\]
非零分支已经稳定。

由\cref{b3:prop:differential-branch-splitting}，
\[
J_i=(J_i:H^\infty)\cap
\bigcap_{j=1}^t\sqrt{[J_i,h_j]}.
\]
初式与分离子都对 \(\mathcal A\) 完全约化。若 \(h_j\) 不是非零常数，把它加入 \(J\) 后即可选出秩严格更小的自约化集；若为非零常数，添入它后得到全环。这是特征集极小秩性质的另一项用途。各链 \(\sqrt{[J_i,h_j]}\) 因而由归纳假设稳定。有限多个退化分支与已经稳定的饱和部分相交，给出原链的稳定性。这里省略特征集秩比较中插入约化多项式的技术归纳；有限性定理及其所需特征集性质见\cite[Section 2.2, Theorem 2.4; Section 3.3]{Hubert2003DifferentialSystems}。

最后，升链条件使逐次选取新的根微分生成元的过程终止。反之，链的并若有有限个根微分生成元，它们已同时落在某一项中，该项以后便不再增长。于是有限基与升链条件两种表述等价。
\end{proofsketch}

\subsection{Rosenfeld–Gröbner 正则微分链与饱和分解}
\label{b3:subsec:E03:02}

局部化把初式与分离子变成可逆元，但这些多项式为零的解仍须另外保存。设 \(S\subset R\) 为有限非零多项式集合，记 \(S^*\) 为其全部幂乘积组成的乘法集，包含空积一。对普通理想 \(J\subseteq R\)，本节采用乘法集饱和记号
\[
J:S^\infty=\{f\in R:hf\in J\;\text{对某个}\;h\in S^*\}.
\]
若 \(h_S=\prod_{s\in S}s\)，则它也等于 \(J:h_S^\infty\)。这是倒置 \(S\) 后理想的收缩；它记录非零条件下的方程后果，其全部零点却仍可能包含 \(h_S=0\) 的特殊点。

\begin{definition}\label{b3:def:regular-differential-system}
设 \((K,\delta)\) 为特征零微分域，\(R=K\{y_1,\ldots,y_n\}\)，并固定导数排序。若有限集 \(\mathcal A\subset R\setminus K\) 的导变量互不相同，且成员两两部分约化，则称 \(\mathcal A\) 为\term{微分三角集}[differential triangular set]。记
\[
\mathcal A_{<v}^{(\delta)}
=\{\delta^jA:A\in\mathcal A,j\ge0,
       \operatorname{leader}(\delta^jA)<v\}.
\]
若 \(S\subset R\setminus\{0\}\) 有限，\(S\) 的成员均对 \(\mathcal A\) 部分约化，且各分离子属于 \(S^*\)，并满足下列相容性条件，则称 \((\mathcal A,S)\) 为\term{正则微分系统}[regular differential system]：对 \(A,B\in\mathcal A\) 的任意共同导变量导数
\(v=\delta^r\operatorname{leader}(A)=\delta^s\operatorname{leader}(B)\)，要求
\[
S_B\delta^rA-S_A\delta^sB
\in\bigl(\mathcal A_{<v}^{(\delta)}\bigr):S^\infty.
\]
这一条件称为相容性（coherence）。系统所定义的正则微分理想为 \([\mathcal A]:S^\infty\)。
\end{definition}

这个定义采用\cite[Definition 4.7]{Hubert2003DifferentialSystems}的形式。在本文的单导子情形，微分三角集的不同成员的导变量来自不同 \(y_i\)，因而没有非平凡的共同导数临界对；偏微分情形则不同。算法的中间列表还未成为三角集时，仍须处理导变量可相互求导得到的归约临界对。正则系统的定义并不要求其初式在全部较早方程模下都是非零因子；这项更强的正则性属于下面的链条件。

正则微分链属于微分三角分解术语。它既不是第15章的正则序列，也不是附录 B 只在有限代数变量中使用的正则链；这里还必须保存导数及其相容条件。

\begin{definition}\label{b3:def:regular-differential-chain}
设 \((K,\delta)\) 为特征零微分域，\(R=K\{y_1,\ldots,y_n\}\)，并固定导数排序。把微分三角集按导变量递增写成 \(\mathcal C=(C_1,\ldots,C_s)\)，记其初式与分离子为 \(I_i,S_i\)，并定义普通代数理想
\[
J_0=(0),\qquad
J_i=(C_1,\ldots,C_i):(I_1\cdots I_i)^\infty
\quad(1\le i\le s).
\]
若各 \(J_i\ne R\)，各 \(I_i\) 在 \(R/J_{i-1}\) 中不是零因子，各 \(S_i\) 在 \(R/J_i\) 中不是零因子，且 \(\mathcal C\) 对其初式与分离子集合满足上述相容性，则称 \(\mathcal C\) 为\term{正则微分链}[regular differential chain]。记
\[
H_{\mathcal C}=\prod_{i=1}^sI_iS_i,
\]
空链时取 \(H_{\mathcal C}=1\)。
\end{definition}

初式的条件使三角方程能逐层作除法，分离子的条件排除这些代数层中的重根。这正是“无重根的正则代数链加上微分相容性”的条件；它与 \(\mathcal C\) 能用伪余式刻画 \([\mathcal C]:H_{\mathcal C}^\infty\) 等价，见\cite[Section 5.1, Theorem 5.2 and Corollary 5.3]{Hubert2003DifferentialSystems}。一般正则系统未必已经给出这样的链，仍可能需要进一步分解。

四类对象的条件及余式意义可以并列如下；成员判据中的理想必须连同饱和条件写出。
\begin{center}
\begin{tabularx}{\linewidth}{@{}>{\raggedright\arraybackslash}p{0.23\linewidth}>{\raggedright\arraybackslash}p{0.33\linewidth}X@{}}
\toprule
对象 & 指定条件 & 余式或成员关系 \\
\midrule
特征集 & 给定微分理想中的极小秩自约化集 & 成员约化为零；反向只立即得到饱和成员关系 \\
正则微分系统 & 三角集、非零条件及相容性 & 部分约化成员可转成代数饱和问题 \\
正则微分链 & 再有初式及分离子的非零因子条件 & 对所定义饱和理想给出完全余式判据 \\
特征展示 & 对正则微分理想作规范化的展示 & 完全零余式与该分支成员关系等价 \\
\bottomrule
\end{tabularx}
\end{center}

\begin{theorem}[Rosenfeld--Gröbner 分解]\label{b3:thm:rosenfeld-groebner}
设 \((K,\delta)\) 为带一个导子的特征零微分域，\(n\ge1\)，\(R=K\{y_1,\ldots,y_n\}\) 的导数排序已给定。设 \(F\subset R\)、\(S\subset R\setminus\{0\}\) 为有限集，\(K\) 中的精确域运算、求导与零判定均可执行。则可计算有限个正则微分链 \(\mathcal C_i\)，使
\[
\sqrt{[F]}:S^\infty
=\bigcap_{i=1}^r[\mathcal C_i]:H_{\mathcal C_i}^{\infty}.
\]
输出可进一步取为特征展示：对任意 \(P\in R\)，
\[
P\in\sqrt{[F]}:S^\infty
\quad\Longleftrightarrow\quad
P\;\text{对每个 \(\mathcal C_i\) 的完全伪余式均为零}.
\]
若左边为全环，输出取空，空交约定为 \(R\)。算法无需在 \(K\) 上作多项式因式分解。
\end{theorem}

这是\cite[Section 6.2, Theorem 60]{BoulierLazardOllivierPetitot2009DifferentialIdeals}的单导子形式。文献中的特征展示具有上述成员判据，并用零维代数 Gröbner 基进一步规范化所选三角列表；具体规范化条件见同文献 Definition 55。特征展示未必已经自约化，但具有正确的饱和成员判据。任意微分理想的 Ritt 特征集则只有“理想成员必约化为零”的方向，不能反过来把所有零余式都算作原理想成员。

这是采用文献的深层算法定理。本节给出输出规格、分支机制和终止量，下面的概要保留“状态与伪约化、初式与分离子分支、严格下降、终端代数细化”四项任务；完整状态更新及其正确性仍取自原文，不能仅据这些段落实现整个算法。

\begin{proofsketch}
算法先把输入分解为正则微分系统，再把各系统细化为提供成员判据的特征展示。前一阶段用状态四元组
\(G=\langle\mathcal A,D,P,S\rangle\) 记录已处理方程、待处理临界对、待处理方程和非零条件；初始状态为 \(\langle\varnothing,\varnothing,F,S\rangle\)。取一个待处理方程或临界对的余式作伪约化，若得到非零多项式 \(q\)，就分别处理其初式、分离子非零的情形及相应因子为零的退化情形。新增方程后同步更新三角列表与临界对。各次更新保存状态不变量；分裂时，各输出分支所定义的根微分理想的交等于原状态的理想。

各次分裂的依据是\cref{b3:prop:differential-branch-splitting}；伪约化保存 \(hP-r\in[\mathcal A]\) 的恒等式。中间状态中，被替换的方程还可能保存在待处理临界对里，不能只把 \(\mathcal A\cup P\) 当作全部方程。完整状态不变量和更新规则见\cite[Sections 5.1--5.3]{BoulierLazardOllivierPetitot2009DifferentialIdeals}。

终止性使用该文 Section 5.4 为状态四元组规定的严格次序：已处理方程列表的秩较小时下降；列表相同时，待处理临界对较少时下降；已处理方程列表与待处理临界对集合均相同时，把某个待处理多项式替换为有限个严格低秩多项式也使状态下降。Lemma 53 通过自约化列表秩的良基性和有限分支树引理证明这个状态次序没有无限下降链；Theorem 28 的证明再按该次序作良基归纳，逐项核对更新后的不变量与下降关系。本概要省去各项更新保持不变量并严格降低状态的技术核对。

终端状态经自约化后成为正则系统。Rosenfeld 引理把部分约化多项式的成员关系化成普通代数饱和的成员关系：
\[
Q\in[\mathcal A]:S^\infty
\quad\Longleftrightarrow\quad
Q\in(\mathcal A):S^\infty
\qquad(Q\;\text{对 \(\mathcal A\) 部分约化}).
\]
Rosenfeld 引理与 Lazard 引理同时给出该微分饱和理想的根性，见\cite[Theorems 23 and 25]{BoulierLazardOllivierPetitot2009DifferentialIdeals}。
Rosenfeld 的论证把求导关系与有限代数计算联系起来。\footnote{罗森菲尔德（Azriel Rosenfeld，1931-2004），美国数学家与计算机科学家，研究微分代数及数字图像分析。}

最后，在包含各方程与非零条件的有限个代数变量中计算饱和的 Gröbner 基。对初式或分离子成为零因子的情形继续分裂，得到正则微分链及特征展示。这个代数细化阶段需要的正则性、根性和规范化论证，见\cite[Section 6.1, Lemmas 57--58 and Theorem 59]{BoulierLazardOllivierPetitot2009DifferentialIdeals}。完全伪余式为零的充要判据来自输出展示的性质，不来自对任意三角列表作伪除法。
\end{proofsketch}

不要求基域多项式因式分解，仍须在有限代数变量中作 Gröbner 基和饱和计算；这两种运算不能混同。输出的饱和理想都是根微分理想，但不要求是素理想；整个分解也不要求互不冗余或唯一。普通 Gröbner 基与这些微分计算的差别集中列在\cref{b3:tab:E-groebner-characteristic-comparison}中。

\begin{table}[htbp]
\centering
\caption{普通 Gröbner 基与微分特征集方法的比较}
\label{b3:tab:E-groebner-characteristic-comparison}
\begin{tabularx}{\textwidth}{@{}>{\raggedright\arraybackslash}p{.17\textwidth}>{\raggedright\arraybackslash}X>{\raggedright\arraybackslash}X@{}}
\toprule
比较项 & 域上的普通多项式理想 & 特征零下的单导子微分理想 \\
\midrule
变量与次序 & 有限变量；单项式序 & 无限导数变量；导数排序及多项式秩 \\
约化中的因子 & 首项系数属于基域，非零时可逆 & 初式与分离子可能是多项式，须记录被倒置的因子 \\
约化证书 & \(P-r\in(G)\) & \(hP-r\in[\mathcal A]\)，其中 \(h\) 为初式与分离子的幂乘积 \\
零余式的含义 & 对 Gröbner 基约化为零当且仅当属于原理想 & 任意特征集的零余式只保证饱和成员关系；正确输出展示才给各分支的充要判据 \\
算法输出 & 原理想的 Gröbner 基 & 根微分理想的有限饱和交分解及各分支的特征展示 \\
\bottomrule
\end{tabularx}
\end{table}

在固定基域的普通 Gröbner 约化中，非零首项系数已可逆；微分约化遇到的初式与分离子却可能随解取零值，所以除法必须连同退化分支处理。表中的微分输出针对根微分理想，不等同于给普通理想找一组有限生成元。

\subsection{消去算例与分母、初式、分离子的退化分支}
\label{b3:subsec:E03:03}

\begin{example}[保留退化分支]\label{b3:exm:differential-singular-branch}
设 \((K,\delta)\) 为特征零微分域，在 \(K\{y\}\) 中，对方程 \(P=(y')^2-y=0\) 求导得到 \(y'(2y''-1)=0\)。在任一微分域解中，或者 \(y'=0\)，从而 \(y=0\)；或者 \(y'\ne0\)，从而 \(y''=\dfrac12\)。若基域含 \(t\) 且 \(t'=1\)，在后一分支令 \(c=2y'-t\)，则 \(c'=0\)，又由原方程得
\[
y=\frac{(t+c)^2}{4}.
\]
故全部微分域解由零解与这族二次式组成，常数 \(c\) 取自所在扩域的常数域。这里 \(y'\ne0\) 是说它不是微分域中的零元素，并不是说某个解析函数的导数在每一点都不为零。若用函数解释这族解，则 \(y'=(t+c)/2\) 在 \(t=-c\) 取零值，却不恒等于零，所以仍属于非零分支。分支区分的是是否恒等满足退化关系，不是将定义域逐点切开；若进一步倒置 \(y'\)，则 \(1/y'=2/(t+c)\) 在解析解释中会在该点出现极点。上述论证逐步检查原式，避免把仅由求导得到的必要条件当作充分条件。
\end{example}

\begin{example}[消去隐藏变量]\label{b3:exm:differential-elimination-system}
设 \((K,\delta)\) 为特征零微分域，在 \(K\{y,z\}\) 中考虑方程组 \(y'=z,z'=yz\)。消去 \(z\) 后得到 \(y''-yy'=0\)，且两者在任意微分域中等价：一方向由求导和代入得到；反方向对满足后一方程的 \(y\) 定义 \(z=y'\)，便验证两条原方程。在微分多项式环层面，有
\[
K\{y,z\}/[y'-z,z'-yz]\cong K\{y\}/[y''-yy'],
\]
同构由 \(z\mapsto y'\) 及恒等的 \(y\) 给出。
\end{example}

第二个例子直接代入 \(z=y'\)，没有倒置初式、分离子或系数，因而给出的是真正的微分商环同构；第一个例子的约去 \(y'\) 则必须另查零解。若系数属于 \(\mathbb Q(t)\)，系数分母在该域中已经可逆；讨论其解析解在某个分母零点附近的行为，又是另一个局部问题。算法报告应分别列明这些不同来源的非零假设。


\begin{example}\label{b3:exm:E-initial-separant}
设 \((K,\delta)\) 为特征零微分域，在 \(K\{y\}\) 中采用 \(y<y'<y''<\cdots\)。对
\[
P=y(y'')^3+(y')^2-1,
\]
微分阶为二，导变量为 \(y''\)，关于导变量的次数为三，初式为 \(y\)，分离子为 \(3y(y'')^2\)。其导数中最高导数 \(y^{(3)}\) 的系数恰为该分离子。
\end{example}
\begin{proof}
把 \(P\) 当作 \(y''\) 的多项式，直接得到初式与次数；对 \(y''\) 作普通偏导得到分离子。指定导子再给出
\[
\delta P=y'(y'')^3+3y(y'')^2y^{(3)}+2y'y''.
\]
最高新变量的系数由此可见。初式与分离子都是非零多项式，求值时却可能为零，因此伪除法使用它们时仍须分别处理相应的非零及退化条件。
\end{proof}

\begin{example}\label{b3:exm:E-branch-certificate}
设 \((K,\delta)\) 为特征零微分域，在 \(K\{y,z\}\) 中取 \(P=yz-y'\)、\(Q=z'-z^2\)。令 \(r=yy''-2(y')^2\)，则有有限等式
\[
y^2Q-r=y\delta P-(yz+2y')P\in[P].
\]
在任意微分域扩张中，原系统 \(P=Q=0\) 恰由两个互不相交分支覆盖：\(y\ne0\) 时为 \(r=0\)、\(z=y'/y\)；\(y=0\) 时须保留 \(z'=z^2\)。
\end{example}
\begin{proof}
展开两边可得 \(y^2z'-y^2z^2-yy''+2(y')^2\)，这验证了乘子证书。在非零分支，\(P=0\) 给出 \(z=y'/y\)，再用分式求导得到
\[
Q=\delta(y'/y)-(y'/y)^2=\dfrac{yy''-2(y')^2}{y^2}.
\]
所以该分支的条件与两条原方程等价。在零分支，\(y\) 及其全部导数都为零，\(P=0\) 自动成立，但仍须满足 \(Q=0\)。原解必属于其中一个分支，反向代入也都满足原方程。若只保留 \(r=0\)，在 \(y=0\) 时会错误地允许任意 \(z\)，因此不能丢掉退化分支的剩余方程。
\end{proof}

% !TeX root = ../../Book3.tex
% 纲要标题：### E.4 线性微分理论与相关专题

\section{线性微分理论与相关专题}
\label{b3:sec:E04}

线性微分方程把求导写成有限维空间上的加法算子，但该算子对标量满足带导子的乘法法则。解矩阵与常数域因此共同决定对称群；加入参数后，还要记录额外导子之间的相容关系。

% 纲要细目（与计划纲要同步）：
% * 线性系统、微分模与Picard–Vessiot扩张；代数闭常数域下的存在唯一性回指第五卷附录E.2，用形式指数计算乘法群算例
% * 微分闭域的existentially closed条件、参数化导子、微分代数群及差分代数的基本对象
% * 完整微分代数几何、Kolchin 维数多项式、算法复杂度及微分模型论不纳入本套书的共同证明链
% #### E.4.1 线性系统、微分模与 Picard–Vessiot 扩张
% #### E.4.2 微分闭域、参数化理论、微分代数群与差分代数
% #### E.4.3 微分代数几何、维数多项式与模型论导读

\subsection{线性系统、微分模与 Picard–Vessiot 扩张}
\label{b3:subsec:E04:01}

\begin{definition}\label{b3:def:differential-module}
设 \((K,\delta)\) 为特征零微分域，\(M\) 为有限维 \(K\) 向量空间。若加法映射 \(\nabla:M\to M\) 满足 \(\nabla(am)=\delta(a)\cdot m+a\cdot\nabla(m)\)，则 \((M,\nabla)\) 称为\term{微分模}[differential module]。满足 \(\nabla(m)=0\) 的向量称为水平向量。
\end{definition}

选定基后可写 \(\nabla(v)=\delta(v)-Av\)，于是水平向量满足线性系统 \(Y'=AY\)。若令 \(Y=PZ\)，\(P\in\operatorname{GL}_n(K)\)，则新系统为 \(Z'=(P^{-1}AP-P^{-1}P')Z\)。额外的 \(P'\) 项来自基随导子变化，故不能照普通矩阵相似变换处理。

\begin{definition}\label{b3:def:picard-vessiot-extension}
设 \((K,\delta)\) 为特征零微分域，\(C=\ker\delta\) 为常数域，\(n\ge1\)，\(A\in M_n(K)\)。设 \(L/K\) 为微分域扩张，延拓的导子仍记为 \(\delta\)，\(C_L=\ker(\delta:L\to L)\)。若 \(L\) 含有矩阵 \(Y\in\operatorname{GL}_n(L)\)，满足 \(Y'=AY\)，且 \(L\) 由 \(Y\) 的各项与 \(\det(Y)^{-1}\) 在 \(K\) 上生成为域，并有 \(C_L=C\)，则称 \(L/K\) 为该系统的\term{Picard--Vessiot 扩张}[Picard--Vessiot extension]。\footnote{韦西奥（Ernest Vessiot，1865-1952），法国数学家，研究微分方程的可积性，发展线性微分方程的 Galois 理论。}与导子交换且固定 \(K\) 的 \(L\) 的域自同构组成其微分 Galois 群。
\end{definition}

设 \((K,\delta)\) 为特征零微分域，且常数域 \(C=\ker\delta\) 代数闭。对于任意 \(n\ge1\) 和 \(A\in M_n(K)\)，系统 \(Y'=AY\) 都有 Picard--Vessiot 扩张，且在固定 \(K\) 并与导子交换的同构意义下唯一。存在唯一性定理见\cite[Propositions 1.20, 1.22]{VanDerPutSinger2003GaloisTheory}；第五卷附录 E.2 将给出这一构造及证明。

\begin{proposition}\label{b3:prop:pv-constant-matrices}
设 \((K,\delta)\) 为特征零微分域，\(C=\ker\delta\)，\(n\ge1\)，\(A\in M_n(K)\)，\(L/K\) 为方程 \(Y'=AY\) 的 Picard--Vessiot 扩张，且选定使 \(L=K(Y_{ij},\det(Y)^{-1})\) 的解矩阵 \(Y\in\operatorname{GL}_n(L)\)。每个与导子交换且固定 \(K\) 的自同构 \(\sigma:L\to L\) 唯一确定 \(C_\sigma\in\operatorname{GL}_n(C)\)，使 \(\sigma(Y)=YC_\sigma\)。映射 \(\sigma\mapsto C_\sigma\) 是单射群同态。
\end{proposition}
\begin{proof}
\(\sigma(Y)'=A\cdot\sigma(Y)\)，而 \((Y^{-1})'=-Y^{-1}A\)，故 \((Y^{-1}\cdot\sigma(Y))'=0\)。无新常数条件使矩阵各项属于 \(C\)。它可逆并唯一。复合时常数矩阵被自同构固定，于是 \(C_{\sigma\tau}=C_\sigma C_\tau\)。若 \(C_\sigma=I\)，则 \(\sigma\) 固定域生成元，故为恒等。
\end{proof}

矩阵表示把微分自同构转化为解空间上的常数线性变换，但哪些常数矩阵真正保持域中的全部代数关系，还须由具体系统决定。第一卷附录 D 给出了基本解及微分自同构的算例；第五卷附录 E 将继续讨论群的代数结构、Galois 对应及其与微分模的关系。

\subsection{微分闭域、参数化理论、微分代数群与差分代数}
\label{b3:subsec:E04:02}

微分方程未必在给定微分域中有解，这促使人们扩大“代数闭”的观念。设 \((K,\delta)\) 为只指定一个导子的特征零微分域。若任意以 \(K\) 为系数、在某个微分扩域可解的有限微分多项式方程与不等式系统，在 \(K\) 中已有解，则称 \((K,\delta)\) 为\term{微分闭域}[differentially closed field]。这里不等式指 \(g\ne0\)，并非序关系。按模型论的说法，\(K\) 在特征零单导子微分域的类中满足 existentially closed 条件：以 \(K\) 中元素为参数的存在性条件，一旦能在某个微分扩域中满足，就已经能在 \(K\) 中满足。条件不仅涉及底层域的多项式，还涉及各阶导数。

\ProseParagraphBreak
若方程还依赖参数，可以同时指定主变量导子 \(\delta\) 和参数导子 \(\partial\)，要求二者交换。\(\delta\) 的常数域对 \(\partial\) 稳定，因为 \(\delta a=0\) 蕴含 \(\delta(\partial a)=\partial(\delta a)=0\)；它自身因而带有参数导子。研究与二者同时交换的自同构时，常数矩阵的系数还会受到 \(\partial\) 方程约束，这就是参数化微分 Galois 理论出现微分代数群的缘由。所谓\term{微分代数群}[differential algebraic group]，在矩阵情形可理解为由微分多项式方程定义、且对乘法和取逆封闭的矩阵集合。例如常数非零元组成 \(\{a:a\ne0,\delta a=0\}\)，它是乘法群的微分代数子群。

\ProseParagraphBreak
\term{差分代数}[difference algebra]则指定环自同态 \(\sigma\)，如 \(f(t)\mapsto f(t+1)\)；它研究 \(y(t+1)\) 与 \(y(t)\) 的关系。导子的 Leibniz 法则与自同态的乘法法则不同，因而差分多项式与微分多项式的理想条件并不相同。

\subsection{微分代数几何、维数多项式与模型论导读}
\label{b3:subsec:E04:03}

微分代数几何用有限阶截断，把无穷多个导数暂时收集在有限生成的代数中。设 \(L=K\langle a_1,\ldots,a_n\rangle\)，定义
\[
h_a(s)=\operatorname{trdeg}_K K(a_i,\delta a_i,\ldots,\delta^sa_i:1\le i\le n).
\]
它统计到第 \(s\) 阶为止还保留多少个普通代数自由参数。若 \(a\) 是一个微分超越元，则 \(h_a(s)=s+1\)；若 \(a'=a\) 且 \(a\) 在 \(K\) 上普通超越，则所有高阶导数等于 \(a\)，故 \(h_a(s)=1\)。两种情形都不能只用“方程有几个未知函数”判断。

\ProseParagraphBreak
若一个数值多项式在充分大的整数 \(s\) 上等于 \(h_a(s)\)，就称它为这一生成元组的\term{微分维数多项式}[differential dimension polynomial]。上面两个例子分别得到 \(s+1\) 与 \(1\)；多个交换导子时则按总微分阶截断。一般存在性定理及其与排序、特征集的关系属于本附录未建立的维数理论。有效计算所涉及的问题可读\cite[Section 1]{Hubert2003DifferentialSystems}。截断过程记录的是给定生成元的导数关系，不能把任意有限阶的解都认作可延拓为全部阶的解。

\ProseParagraphBreak
模型论把这些问题写成微分域上的公式：方程表达等式，不等条件表达可逆性，量词表达存在新的解或参数。微分闭域提供研究解集的一种环境，但它不自动指定解析函数，也没有给出收敛区域。前文已经从方程构造微分理想，在消去中保留非零条件，并由线性系统得到常数矩阵作用；更一般的几何与模型论需要另作系统学习。



\begin{chaptersummary}
指定导子改变了允许的理想与映射：求值必须同时规定全部导数的像，商环必须由微分理想形成。微分多项式环虽然有无穷多个普通变量，根微分理想仍有 Ritt--Raudenbush 有限性；其特征集证明先稳定非零饱和分支，再按低秩退化分支归纳。Rosenfeld--Gröbner 方法还须分别控制待约化方程、相容对、初式和分离子，有限性本身不能替代这些检查。

\ProseParagraphBreak
计算的代表任务是在保留解集的前提下消去隐藏变量。非线性情形需要三角分支和饱和；线性情形则由微分模与基本解矩阵组织。两者都属于代数理论，解析解的收敛、定义域与奇点延拓还须另有论证。Kähler 微分、普通导子和 DG 微分之间的区别，也应在进入更深理论后继续保留。
\end{chaptersummary}

\BookExercises

\begin{exercise}\label{b3:ex:E-derivation-constants}
设 \(k\) 为特征零域，在 \(K=k(t)\) 上指定 \(\delta|_k=0\)、\(\delta t=t\)。求常数域以及全部与 \(\delta\) 交换的 \(k\) 自同构。说明普通域自同构 \(t\mapsto t+1\) 为什么不属于这个群。
\end{exercise}
\begin{solution}
导子为 \(t\,\mathrm{d}/\mathrm{d}t\)，因 \(t\ne0\)，常数域与通常求导相同，为 \(k\)。若 \(\sigma\) 与导子交换，令 \(r=\sigma(t)\)，则 \(\delta r=r\)，于是 \(\delta(r/t)=0\)。因此 \(r=ct\)、\(c\in k^\times\)。反过来每个这种代换都定义域自同构，在 \(k,t\) 上与导子交换，故在整个域上交换。群为 \(k^\times\)。对平移有 \(\delta(t+1)=t\)，而 \(\sigma(\delta t)=t+1\)，两者不同。
\end{solution}

\begin{exercise}\label{b3:ex:E-polynomial-differential-ideals}
设 \(K\) 为特征零域，在 \(K[t]\) 上指定 \(\delta|_K=0\)、\(\delta t=1\)。证明该环只有零理想和全环两个微分理想。
\end{exercise}
\begin{solution}
若微分理想含非零多项式，取次数最小者 \(p\)。若次数正，则 \(p'\) 非零且次数更小，也在理想中，矛盾。因此理想含非零常数，为全环。
\end{solution}

\begin{exercise}\label{b3:ex:E-characteristic-p-warning}
令 \(p\) 为素数，在 \(A=\mathbb F_p[t]\) 及其分式域上取 \(\delta t=1\)。求分式域的全部常数，分类 \(A\) 的非零微分理想，并确定其中哪些理想的根仍是微分理想。
\end{exercise}
\begin{solution}
若既约非零分式 \(a/b\) 被导子零化，\(a'b=ab'\) 及互素性给出 \(a\mid a'\)。次数比较使 \(a'=0\)，进而 \(b'=0\)，所以 \(a,b\in\mathbb F_p[t^p]\)。反向求导为零，故常数域恰为 \(\mathbb F_p(t^p)\)。非零理想可写成 \((g)\)；它微分稳定当且仅当 \(g'\in(g)\)，同样由次数得 \(g'=0\)，即 \(g\in\mathbb F_p[t^p]\)。因为系数在 \(\mathbb F_p\)，非常数的这种 \(g\) 是一个多项式的 \(p\) 次幂，其根理想的非恒定无重因子生成元不可能导数为零。因此非零微分理想中仅全环的根仍微分稳定；零理想也为根微分理想。
\end{solution}

\begin{exercise}\label{b3:ex:E-differential-polynomial-map}
计算 \(K\{y\}/[y'-y]\) 的底层环，并说明它不是只含常数的微分环。
\end{exercise}
\begin{solution}
所有关系是 \(y^{(j+1)}-y^{(j)}\)，故各高阶变量都等于 \(y\)，商的底层环为 \(K[y]\)。令 \(\delta y=y\) 可定义逆向微分同态，验证没有额外关系。除非 \(y=0\)，生成元的导数非零，所以它并非零导子结构。
\end{solution}

\begin{exercise}\label{b3:ex:E-root-differential-ideal}
在特征零微分域上证明 \(\sqrt{[y^2]}=[y]\)。
\end{exercise}
\begin{solution}
\([y]\) 由全部 \(y^{(j)}\) 生成，其商为 \(K\)，所以是根微分理想并包含 \([y^2]\)。反之 \(y^2\in[y^2]\)，故 \(y\in\sqrt{[y^2]}\)；正文引理保证后者是微分理想，从而包含 \([y]\)。
\end{solution}

\begin{exercise}\label{b3:ex:E-initial-separant}
设 \((K,\delta)\) 为特征零微分域，\(P=z(y')^2+yz'+1\in K\{y,z\}\)。分别采用按阶排序 \(y<z<y'<z'<\cdots\) 和消去排序 \(z<z'<\cdots<y<y'<\cdots\)，计算导变量、关于导变量的次数、初式和分离子，并核对 \(\delta P\) 的最高新变量及其系数。
\end{exercise}
\begin{solution}
按阶排序的导变量为 \(z'\)，次数为一，初式与分离子都为 \(y\)。消去排序的导变量为 \(y'\)，次数为二，初式为 \(z\)，分离子为 \(2zy'\)。展开
\[
\delta P=z'(y')^2+2zy'y''+y'z'+yz''.
\]
第一种排序的最高变量为 \(z''\)，系数是 \(y\)；第二种排序的最高变量为 \(y''\)，系数是 \(2zy'\)。微分阶仍为一，普通总次数仍为三，但导变量与伪除法需要的因子随排序改变。
\end{solution}

\begin{exercise}\label{b3:ex:E-branch-certificate}
设 \((K,\delta)\) 为特征零微分域，在 \(K\{y,z\}\) 中取 \(P=yz-y'-1\)、\(Q=z'\)。构造消去 \(z\) 的乘子证书，并证明原系统在微分域扩张中与 \(y\ne0\)、\(yy''-(y')^2-y'=0\)、\(z=(y'+1)/y\) 等价。说明只保留消去方程会多出什么解。
\end{exercise}
\begin{solution}
记 \(r=yy''-(y')^2-y'\)，直接展开给出
\[
y^2Q-r=y\delta P-y'P\in[P].
\]
若 \(y=0\)，它的所有导数均为零，\(P=-1\)，故原系统没有这个分支。对 \(y\ne0\)，第一式给出 \(z=(y'+1)/y\)，且
\[
z'=\dfrac{yy''-(y'+1)y'}{y^2}=\dfrac r{y^2}.
\]
于是两组条件正反等价。消去方程单独允许 \(y=0\)，这个解不能提升到原系统；乘子条件不能从输出中删去。
\end{solution}

\begin{exercise}\label{b3:ex:E-division-loses-solution}
在微分域中分析 \(yy'=0\)，指出直接约去 \(y\) 的逻辑问题，并说明最终解集是什么。
\end{exercise}
\begin{solution}
域无零因子，故 \(y=0\) 或 \(y'=0\)。第一种也是第二种的解，最终解集为全部常数。直接约去 \(y\) 仅在 \(y\ne0\) 分支合法；虽然本例遗漏的零解最后也满足 \(y'=0\)，仍应通过原方程确认，不能据此省略一般分支检查。
\end{solution}

\begin{exercise}\label{b3:ex:E-linear-elimination}
证明系统 \(u'=v,\ v'=-u\) 等价于 \(u''+u=0\) 连同 \(v=u'\)。再证明 \(u^2+v^2\) 为常数。
\end{exercise}
\begin{solution}
第一式求导并用第二式得到 \(u''=-u\)；反过来定义 \(v=u'\) 即满足原系统。再算 \(\delta(u^2+v^2)=2uv+2v(-u)=0\)。此常数不必属于最初的底域，取决于所用微分扩域的常数域。
\end{solution}

\begin{exercise}\label{b3:ex:E-riccati-pole}
设 \(y\ne0\) 满足 \(y'=y^2\)，基域含 \(t\) 且 \(t'=1\)。证明 \(-1/y-t\) 为常数，并解释零解为何需另列。
\end{exercise}
\begin{solution}
\(\delta(-1/y)=y'/y^2=1\)，故 \(c=-1/y-t\) 为常数，解写成 \(y=-1/(t+c)\)。计算使用 \(y^{-1}\)，所以仅覆盖非零解。\(y=0\) 直接满足原方程，必须另列。
\end{solution}

\begin{exercise}\label{b3:ex:E-module-basis-change}
对一维系统 \(y'=ay\)，令 \(y=pz\)，\(p\in K^\times\)。求新系数，并说明何时能变成 \(z'=0\)。
\end{exercise}
\begin{solution}
\(p'z+pz'=apz\)，所以 \(z'=(a-p'/p)z\)。当且仅当存在 \(p\in K^\times\) 满足 \(p'/p=a\) 时可把系数消去。这是微分域中一个方程的可解性问题，不能仅凭 \(a\) 是标量就宣布系统平凡。
\end{solution}

\begin{exercise}\label{b3:ex:E-constant-solution-space}
设 \(L/K\) 常数域同为 \(C\)，\(Y\in\operatorname{GL}_n(L)\) 满足 \(Y'=AY\)。证明 \(L^n\) 内该系统的全部解为 \(Yc\)，\(c\in C^n\)。
\end{exercise}
\begin{solution}
对任一解 \(v\)，计算 \((Y^{-1}v)'=-Y^{-1}Av+Y^{-1}Av=0\)，故 \(Y^{-1}v\in C^n\)。反方向求导 \(Yc\) 即验证。可逆性保证 \(c\) 唯一，故解空间作为 \(C\) 向量空间维数为 \(n\)。
\end{solution}

\begin{exercise}\label{b3:ex:E-pv-exponential-group}
在 \(K=\mathbb C(t)\) 上取通常导子，令 \(u=\sum_{j\ge0}t^j/j!\in\mathbb C(\!(t)\!)\)、\(L=K(u)\)。证明 \(Y=\operatorname{diag}(u,u^2)\) 给出系统 \(Y'=\operatorname{diag}(1,2)Y\) 的 Picard--Vessiot 扩张，并求其常数矩阵群。若改用基本解矩阵 \(\widetilde Y=YP\)，其中 \(P=\left(\begin{smallmatrix}1&1\\0&1\end{smallmatrix}\right)\)，求相应的新矩阵表示。
\end{exercise}
\begin{solution}
由\cref{b3:prop:E-formal-exponential-pv}，\(u\) 超越且 \(L\) 没有新常数。矩阵 \(Y\) 可逆，满足所给系统，它的各项生成 \(L\)，故 PV 条件全部满足。每个微分自同构为 \(u\mapsto cu\)、\(c\in\mathbb C^\times\)，必须同时把 \(u^2\) 送到 \(c^2u^2\)，所以常数矩阵恰为 \(\operatorname{diag}(c,c^2)\)，不能独立选择两个倍数。\(P\) 为常数矩阵，\(\widetilde Y\) 仍是同一系统的基本解矩阵；新表示为
\[
P^{-1}\operatorname{diag}(c,c^2)P
=\begin{pmatrix}c&c-c^2\\0&c^2\end{pmatrix}.
\]
这个共轭改变表示的基，没有改变扩域或自同构群。
\end{solution}

\begin{exercise}\label{b3:ex:E-commuting-derivations}
在 \(K[x,y]\) 上取 \(\delta_1=\partial/\partial x\)、\(\delta_2=x\partial/\partial y\)。证明二者不交换，并说明为何不能直接套用本附录的多导子基定理。
\end{exercise}
\begin{solution}
\(\delta_1\delta_2(y)=\delta_1(x)=1\)，而 \(\delta_2\delta_1(y)=0\)。所以混合导数不能只用多重指标 \((a,b)\) 表示而忽略次序。正文陈述的是单导子 Ritt--Raudenbush 定理；所提到的多导子一般形式要求导子两两交换。此例不符合后一假设，混合导数还须按实际复合次序记录。
\end{solution}
